\documentclass{tietreport}

% Import images
\usepackage{graphicx}
\usepackage{booktabs}
\usepackage{float}

% Bibliography configuration
\usepackage[
  date=year,
  style=alphabetic,
  backend=biber,
  sorting=nty,
  maxnames=5,
  minnames=3,
  maxalphanames=4,
  minalphanames=3,
  backref=true,
  doi=false,
  isbn=false,
  url=false,
  eprint=false
]{biblatex}
\addbibresource{references.bib}

% Give TeX room to resolve awkward lines rather than
% overflowing the margin.
\setlength{\emergencystretch}{3em}

%% ==================================================
%% CUSTOMIZE THESE FIELDS FOR YOUR PROJECT
%% ==================================================

% Course/Lab header
\TitlePageHeader{%
  UCS 503 Software Engineering Project Proposal%
}

% Main project title
\ProjectTitle{%
  CampusCycle%
}

% Subtitle or project description
\ProjectSubTitle{%
  Smart Campus Bicycle Sharing and
  Mobility Management System%
}

% Student information
\author{%
  \texttt{102XXXXXXX} Guneet Singh Nagi \\
  \texttt{102XXXXXXX} Prabhnoor Oberoi%
}

% Group number, degree, year, program
\TitlePageSubText{%
  Group: \texttt{XXXX} BE Third Year, CSE%
}

% Faculty advisor/guide name
\AdvisorName{%
  Dr. Faculty Member Name%
}

% Evaluation period and department info
\TitlePageFooterText{{%
    \large Project Proposal \\[0.2cm]
    \bfseries Computer Science and Engineering
    Department%
  } \\[0.15cm]
  Thapar Institute of Engineering and Technology,
  Patiala%
}

%% ==================================================

\begin{document}

% Generate title page
\maketitle

% Table of contents (auto-generated from chapters)
\tableofcontents

% ===== CHAPTER 1: INTRODUCTION =====

\chapter{Introduction}

\section{Background and Context}

CampusCycle is a smart bicycle-sharing platform designed
specifically for university campuses. The system allows
students to locate, reserve, unlock, use and return
shared bicycles through a web application.
Administrators can manage bicycles, stations, rentals,
maintenance and usage analytics.

The project extends beyond basic rental functionality by
incorporating QR-based bicycle identification,
geofencing, maintenance reporting, usage analytics and
demand prediction. These features aim to make
short-distance transportation within a large campus more
convenient, affordable, sustainable and manageable.

Shared-bicycle schemes are well established as a
solution to the last-mile problem and have been deployed
in several hundred cities
worldwide~\cite{shaheen2010}. The research literature
on these schemes consistently identifies two factors as
decisive for whether a scheme is adopted or abandoned:
reliable information about where bicycles are available,
and the condition of the fleet~\cite{fishman2016}. Both
are information problems rather than hardware problems,
which is what makes a campus scheme a suitable software
engineering project.

\subsection{Problem Statement}

Large university campuses often require students to
travel considerable distances between hostels, academic
blocks, libraries, cafeterias, sports facilities and
administrative buildings. Although bicycles can solve
this last-mile mobility problem, a shared bicycle system
requires efficient tracking and management.

A conventional rental system does not adequately address
problems such as uneven bicycle distribution,
unauthorised use outside the campus, maintenance issues,
and difficulty in predicting where bicycles will be
needed. Studies of operating schemes report that
redistributing bicycles between stations is one of the
largest recurring operational
costs~\cite{faghihimani2017}, and that cost is driven by
decisions taken without reliable demand information.

Therefore, the proposed system aims to develop a
centralised smart campus bicycle-sharing platform that
manages the bicycle fleet while using location and
historical usage data to improve availability and
operational efficiency.

\section{Project Objectives}

The main objectives of CampusCycle are:

\begin{enumerate}
  \item Provide students with a simple platform to find
    and rent campus bicycles.
  \item Implement QR-based identification for individual
    bicycles.
  \item Display bicycle and docking-station availability
    on an interactive campus map.
  \item Maintain complete rental and return records.
  \item Use geofencing to detect bicycles leaving the
    permitted campus area.
  \item Allow students to report bicycle faults and
    maintenance problems.
  \item Provide administrators with fleet, usage and
    maintenance analytics.
  \item Analyse historical rental data to predict
    bicycle demand at different stations.
  \item Recommend redistribution of bicycles between
    stations based on predicted demand.
\end{enumerate}

% ===== CHAPTER 2: PROPOSED SYSTEM =====

\chapter{Proposed System}

The system will contain two major interfaces: a student
application and an administrator dashboard.

\section{Student Application}

Students will be able to:

\begin{itemize}
  \item Register and log in using their university
    account.
  \item View nearby bicycle stations and available
    bicycles.
  \item Scan a QR code to identify a bicycle.
  \item Reserve and start a bicycle rental.
  \item View the active ride and rental duration.
  \item Return the bicycle at an approved station.
  \item View previous rides and distance travelled.
  \item Report mechanical or safety issues.
  \item View personal mobility statistics and
    sustainability contributions.
\end{itemize}

\section{Administrator Dashboard}

Administrators will be able to:

\begin{itemize}
  \item Add, remove and update bicycles and stations.
  \item Monitor active rentals.
  \item Track bicycle availability and locations.
  \item View maintenance reports and bicycle status.
  \item Detect geofence violations.
  \item Analyse station-wise and time-wise usage.
  \item View predicted demand for upcoming periods.
  \item Receive recommendations for bicycle
    redistribution.
\end{itemize}

\section{Key Features}

\subsection{QR-Based Bicycle Identification}

Each bicycle will have a unique QR code. Scanning the
code allows the backend to identify the bicycle and
verify whether it is available, reserved, currently
rented or under maintenance.

\subsection{Campus Geofencing}

A virtual boundary will be defined around the campus. In
an extended implementation with GPS-enabled hardware,
the system can detect when a bicycle crosses this
boundary and generate an alert for the administrator.
For the software prototype, location data can be
simulated or obtained from a mobile device.

\subsection{Interactive Campus Map}

The application will display stations and their current
availability. Stations can be represented using status
indicators such as available, low availability and
unavailable.

\subsection{Maintenance Management}

Students can report issues such as damaged brakes, flat
tyres or faulty lights. Repeated reports can increase the
priority of a maintenance request, and bicycles marked
unsafe can automatically become unavailable for new
rentals.

\subsection{Demand Prediction}

Historical rental information will be analysed using
machine-learning techniques. Features may include
station, time of day, day of week, previous demand and
number of active rentals. A regression or tree-based
model will predict future demand.

\subsection{Smart Redistribution}

Predicted demand will be compared with current station
capacity. The system can recommend moving bicycles from
stations with excess supply to stations where demand is
expected to be high.

% ===== CHAPTER 3: ARCHITECTURE & TECHNOLOGY =====

\chapter{System Architecture and Technology}

\section{Proposed Architecture}

The proposed architecture consists of:

\begin{enumerate}
  \item \textbf{Frontend:} React.js~\cite{reactdocs}
    with Tailwind CSS for the student and administrator
    interfaces.
  \item \textbf{Backend:} Node.js and
    Express.js~\cite{expressjs} for authentication,
    rentals, bicycle management, maintenance and REST
    APIs.
  \item \textbf{Database:} MySQL for users, bicycles,
    stations, rentals, locations, maintenance records
    and analytics data.
  \item \textbf{Machine Learning Module:} Python with
    Scikit-learn~\cite{scikitlearn} for demand
    prediction and redistribution recommendations.
  \item \textbf{Mapping:} Leaflet~\cite{leaflet} with
    OpenStreetMap~\cite{openstreetmap}, or another
    suitable mapping service, for the campus map.
  \item \textbf{Optional Hardware Extension:} ESP32 with
    GPS and a lock mechanism for demonstrating real-time
    bicycle tracking.
\end{enumerate}

\section{Technology Stack}

\begin{table}[H]
\centering
\begin{tabular}{@{}ll@{}}
\toprule
\textbf{Component} & \textbf{Technology} \\
\midrule
Frontend          & React.js, Tailwind CSS \\
Backend           & Node.js, Express.js \\
Database          & MySQL \\
Machine Learning  & Python, Scikit-learn, Pandas \\
Maps              & Leaflet / OpenStreetMap \\
Authentication    & JWT~\cite{rfc7519} / University
                    Authentication \\
Version Control   & Git and GitHub \\
Optional Hardware & ESP32, GPS module, electronic lock \\
\bottomrule
\end{tabular}
\caption{Proposed technology stack.}
\label{tab:stack}
\end{table}

\section{Project Scope}

The minimum viable product will include authentication,
bicycle and station management, QR-based identification,
rental and return workflows, campus map visualisation,
maintenance reporting, and an administrator dashboard.

Advanced features such as real-time GPS tracking,
electronic locks, demand prediction and smart
redistribution will be implemented as higher-level
modules depending on available development time and
hardware availability.

% ===== CHAPTER 4: OUTCOME & CONCLUSION =====

\chapter{Expected Outcome and Conclusion}

\section{Expected Outcome}

At the end of the project, a working prototype of
CampusCycle will allow a student to discover an
available bicycle, scan its QR code, start a rental,
view the ride status and return the bicycle at an
approved station. Administrators will have a dashboard
for fleet and maintenance management.

The system will additionally demonstrate analytics and a
machine-learning model capable of predicting
station-level bicycle demand. The prototype will provide
a foundation that could later be extended to a real
campus deployment with GPS trackers, electronic locks,
payment integration and mobile applications.

\section{Conclusion}

CampusCycle combines web development, databases,
geospatial concepts, machine learning and optional IoT
hardware into a practical campus mobility solution. The
project addresses a real transportation problem while
remaining feasible as a two-person semester project. Its
modular design also allows the basic bicycle-sharing
system to be developed first and advanced intelligence
and hardware features to be added incrementally.

% ===== BIBLIOGRAPHY =====

% Print bibliography with heading in TOC
\printbibliography[heading=bibintoc]

\end{document}
